\documentclass[11pt,letterpaper]{article}
\usepackage[margin=1in]{geometry}
\usepackage{amsmath,amssymb}
\usepackage{fancyhdr}
\usepackage{xcolor}
\usepackage{enumitem}
\usepackage[framemethod=tikz]{mdframed}
\usepackage[hidelinks]{hyperref}
\pagestyle{fancy}
\fancyhf{}
\lhead{\small AntsPhysicsLessons.com}
\rhead{\small Electricity \& Magnetism}
\cfoot{\small\thepage}
\renewcommand{\headrulewidth}{0.4pt}
\setlength{\parindent}{0pt}
\definecolor{keybg}{RGB}{240,244,250}
\newmdenv[backgroundcolor=keybg,linewidth=0pt,innerleftmargin=8pt,innerrightmargin=8pt,innertopmargin=6pt,innerbottommargin=6pt,skipabove=6pt]{answer}
\begin{document}
{\small\textsc{Unit 1 -- Electric Charge, Force, and Field}}\\[2pt]{\Large\bfseries 1.1 Coulomb's Law and Superposition}\\[8pt]\textbf{Answer key}\\[4pt]\hrule\vspace{10pt}{\small\itshape Placeholder worksheet for the site prototype. Free for classroom use.}\vspace{6pt}
\begin{enumerate}[leftmargin=*,itemsep=10pt]
\item Two point charges, $q_1 = +3.0\,\mu\text{C}$ and $q_2 = -5.0\,\mu\text{C}$, are $0.20$~m apart. Find the magnitude of the force each exerts on the other and state whether it is attractive or repulsive.
\begin{answer}$F = \dfrac{k|q_1 q_2|}{r^2} = \dfrac{(8.99\times10^{9})(3.0\times10^{-6})(5.0\times10^{-6})}{(0.20)^2} = 3.4\ \text{N}$. Opposite signs, so the force is \textbf{attractive}.\end{answer}
\item Three charges lie on the $x$-axis: $+2.0\,\mu\text{C}$ at $x=0$, $-2.0\,\mu\text{C}$ at $x=0.10$~m, and $+4.0\,\mu\text{C}$ at $x=0.30$~m. Find the net force (magnitude and direction) on the $+4.0\,\mu\text{C}$ charge.
\begin{answer}From the $+2.0\,\mu$C charge (repulsive, $+x$): $F_1 = \dfrac{k(2.0\times10^{-6})(4.0\times10^{-6})}{(0.30)^2} = 0.80$~N. From the $-2.0\,\mu$C charge (attractive, $-x$): $F_2 = \dfrac{k(2.0\times10^{-6})(4.0\times10^{-6})}{(0.20)^2} = 1.80$~N. Net: $F = 0.80 - 1.80 = -1.0$~N, i.e.\ \textbf{1.0~N toward $-x$}.\end{answer}
\item Two identical small spheres carry charges $q$ and $3q$ and repel with force $F$. They are touched together and separated to the same distance. Find the new force in terms of $F$.
\begin{answer}After contact each carries $2q$. The product of charges goes from $3q^2$ to $4q^2$, so the new force is $\tfrac{4}{3}F$.\end{answer}
\end{enumerate}
\end{document}